% THE Q-MODEL (TQM) — Executive Summary
% Publication-quality condensed overview for Zenodo readers.
% Source: TQM_v1_0_Monograph_Final.tex (no new physics, no changed conclusions).
% Build: pdflatex TQM_v1_0_Executive_Summary.tex (twice)

\documentclass[10pt,a4paper]{article}
\usepackage[utf8]{inputenc}
\usepackage[T1]{fontenc}
\usepackage{amsmath,amssymb}
\usepackage[margin=2.2cm]{geometry}
\usepackage{booktabs}
\usepackage{array}
\usepackage{lmodern}
\usepackage[hidelinks]{hyperref}
\usepackage{xcolor}

\definecolor{caveat}{HTML}{7A2E2E}

\newcommand{\derived}{\textbf{DERIVED}}
\newcommand{\realunder}{\textbf{REAL-UNDERIVED}}
\newcommand{\drawn}{\textbf{DRAWN}}

\setlength{\parindent}{0pt}
\setlength{\parskip}{4pt plus 1pt}

\title{\bfseries THE Q-MODEL (TQM) — Executive Summary\\[2mm]
\large A Condensed Overview for Zenodo Readers}
\author{Fabrice Wieser\\[1mm]
\small Independent Researcher and Full Stack Software Engineer\\[0.5mm]
\small \url{https://github.com/MagusDraconis/AT}}
\date{Version 1.0 (final) \quad·\quad 15 August 2026}

\begin{document}
\maketitle

\begin{center}
\fcolorbox{caveat}{white}{%
\begin{minipage}{0.94\textwidth}
\color{caveat}
\small
\textbf{Publication status: READY AS A FOUNDATION MONOGRAPH --- READY AS A
RESEARCH-PROGRAM PAPER --- NOT READY AS A DERIVATION PAPER.}\\[1.5mm]
This summary is a condensation of \emph{THE Q-MODEL --- From Q to Cosmology}
(\texttt{TQM\_v1\_0\_Monograph\_Final.tex}). It introduces no new physics, changes no
conclusion, and adds no new derivation; it only summarizes material already present in the
full monograph.
\end{minipage}}
\end{center}

\vspace{2mm}

%=============================================================================
\section{Executive Overview}
%=============================================================================

\textbf{What is TQM.} THE Q-MODEL (TQM) is a theory of \emph{structure} and \emph{content}
that compresses observable physics to a minimal primitive set. It is a foundation
monograph and a research-program record, not a claimed derivation of the Standard Model
from first principles.

\textbf{Core claim.} The single organizing claim is the \emph{structure/content split}:

\begin{center}
\emph{Structure is derivable; content is realized.}
\end{center}

\textbf{Structure versus content.} The \emph{form} of observable structure --- topology,
symmetry, distribution shape, lower bounds --- is provable from a minimal primitive set.
The \emph{content} --- specific masses, couplings, multiplicities, angles --- is a draw
from a contingent ensemble and is therefore \emph{classified} rather than \emph{derived}.
The theory answers \emph{what exists and why} for structure, and honestly declines to
derive the dimensionless numbers that fix \emph{how much}.

Three statuses govern every result: \derived\ (computable from the primitives by a
theorem), \realunder\ (real, precise, predictive structure whose origin is not computable
from the primitives), and \drawn\ (a coincidental draw of the abundance law).

%=============================================================================
\section{Motivation}
%=============================================================================

\textbf{Why TQM was created.} TQM began from a single question: are matter, quantum
behaviour, and gravitation \emph{fundamental}, or do they emerge from something more
basic? The guiding hypothesis is that \emph{matter is not fundamental} --- that a particle
is a dynamically stabilized wave structure, and that synchronization and resonance are
more fundamental than particles. The research path was fixed early and never reversed:
temporal oscillators $\to$ synchronization $\to$ resonance clusters $\to$
self-stabilizing wave structures $\to$ proto-particles $\to$ quantum behaviour $\to$
gravity $\to$ cosmology. Four standing rules disciplined the program: no cosmology first,
no gravity first, no particle assumptions, and emergence first, with physics before
interpretation.

\textbf{The TRM $\to$ TQM transition.} The immediate predecessor was the \emph{Temporal
Resonance Model} (TRM), a phase-lattice framework with nine modules (Time Field, Radial
Acceleration Relation, Frame Dragging, Memory Channel, Theta Chain, $m=3$ Closure, Quantum
Engine, Temporal Drift, and a Unified Action roadmap). The TRM legacy audit found that TRM
carried \emph{zero candidate physics} --- no module supplied a genuinely new testable
observable. Three modules were absorbed (the Time Field became phase-gradient gravity; the
RAR became $g_\dagger=cH_0/2\pi$; Frame Dragging was recognized as GR gravitomagnetism),
two were rejected (tired-light cosmology; a non-unitary Quantum Engine), and three survived
as candidate mathematics only. The transition therefore produced a \emph{clean
reconstruction}: the new theory re-derived from the fewest possible ingredients exactly the
structure TRM had assembled piecemeal, and discarded the rest.

\textbf{Development timeline in brief.} The program ran in five eras. (1) \emph{Early
oscillator experiments} (TQM-001--008) established that synchronization emerges
spontaneously, that random matrices and static topology are insufficient, and that a sharp
critical resonance density ($\rho_c\approx0.09$) marks the onset of persistent coherent
modes --- while a reproducibility study collapsed an apparent five-family diversity to a
single universal mode (F4). (2) \emph{Alternative foundations} (ResearchX, X001--X034)
proved that reversibility plus self-consistency is minimally sufficient for unitary quantum
mechanics, and that quantum reality is the unique global maximum of finite complexity.
(3) \emph{Spacetime emergence} (X040--X046) derived time as a partial order, gravity as
causal structure, and $3{+}1$ from complexity. (4) \emph{Matter emergence} (X047--X060)
derived particles as topological defects and $U(1)$ from the circle. (5) \emph{Parameter
compression and closure} (Phases 140--159) reduced the theory to two primitives plus one
number and closed the residual questions.

%=============================================================================
\section{The Four Primitives}
%=============================================================================

The theory admits exactly four primitives and no others.

\begin{enumerate}
\item \textbf{$Q$ --- the individuation principle.} The irreducible act by which a discrete
event comes to be individuated from nothing. $Q$ underwrites the derivation of structure.
Its formal status is \emph{partially formalized}: object, state space, and operations
exist; the measure and the action functional remain postulates.
\item \textbf{Random Actualization (assumption A-03).} Genuine ontological chance: given
$Q$, the realized event locations are random within the causal structure. This underwrites
content. Its random variable and generator are formalized; the probability space and
ensemble measure are only partially formalized.
\item \textbf{$(\ell,\tau,\hbar)$ --- the irreducible physical triple.} A spacetime scale
$\ell$, a clock $\tau$, and an action $\hbar$ (unit conventions). These fix units; they
are not free parameters.
\item \textbf{$M^2$ --- the single continuous nonlinearity parameter.} The one continuous
number, the nonlinearity regime of the dynamics, pinned by the derivation hierarchy to
$M^2\approx 5$.
\end{enumerate}

No gauge-group, multiplicity, or hidden-dimension primitive is permitted. The primitives
meet in $Q$: an \emph{ontology} layer underwrites structure, and a \emph{dynamical} layer
(the graph Laplacian $L_Q$, the Hilbert space, and the Schr\"odinger equation) underwrites
the dynamics.

%=============================================================================
\section{Main Results}
%=============================================================================

\subsection{Derived}

\begin{itemize}
\item \textbf{$U(1)$ (confidence 0.95).} Phase lives on the circle $S^1$; its isometry
group is $U(1)$, and the winding yields integer charge:
\begin{equation}
\mathrm{Aut}(S^1)=U(1),\qquad \pi_1(S^1)=\mathbb Z .
\end{equation}
$U(1)$ is not chosen; it is the symmetry of phase itself. (Structure only: the Maxwell
action is still borrowed.)
\item \textbf{$3{+}1$ dimensionality (confidence 0.85).} $d=3{+}1$ is the unique value
satisfying simultaneously Bertrand's theorem, Gauss's law, knot stability, Huygens's
principle, and the two GR polarizations. It is a \emph{window-intersection} argument, not a
variational theorem.
\item \textbf{Schr\"odinger dynamics.} Conserved norm forces the generator to be
anti-Hermitian; the simplest anti-Hermitian object on the graph is $J\otimes L_Q$ with
$J^2=-I$, giving
\begin{equation}
i\,\partial_t\psi = L_Q\,\psi ,
\end{equation}
with unitary evolution $\psi(t)=e^{-iL_Qt}\psi(0)$. Quantum mechanics is produced from a
conservation requirement, not postulated.
\item \textbf{Abundance-law structure.} A multiplicative actualization cascade plus the
central-limit theorem in log-space yields a log-normal distribution \emph{form} --- a
theorem --- while the realized values $\mu,\sigma$ are contingent content.
\end{itemize}

\subsection{Real-Underived}

\begin{itemize}
\item \textbf{The Koide relation.} The charged-lepton pole masses satisfy
\begin{equation}
Q=\frac{\sum m}{(\sum\sqrt m)^2}=2/3\qquad(\theta=45^\circ),
\end{equation}
to $\sim10^{-5}$, predictive (1981 prediction of $m_\tau$, confirmed 1992), and RG-stable.
Its origin is \emph{underivable} (no-go T-08): seven routes --- symmetry, topology,
attractors, information geometry, group theory, $m=3$ closure, moduli --- all failed or
were blocked. It is the canonical \realunder\ structure.
\item \textbf{$SU(2)$ (confidence 0.70).} Emergent: the binary winding pair $\{n=\pm1\}$
gives the minimal doublet, the Bloch sphere gives $SO(3)=SU(2)/Z_2$, and the spinor double
cover gives $SU(2)$; the lift from $Z_2$ to continuous $SU(2)$ is not itself derived.
\item \textbf{$SU(3)$ structure (confidence 0.10).} The
$\mathrm{Aut}(C^n/S_n)\supseteq SU(n)$ argument derives the group \emph{structure} but not
the \emph{count}, and the 8-gluon algebra is borrowed (A-07).
\end{itemize}

\subsection{Drawn}

\begin{itemize}
\item \textbf{Couplings} ($\alpha$, $\alpha_s$, $\theta_W$, Yukawas). Dimensionless numbers
cannot be assembled from the dimensionful triple $(\ell,\tau,\hbar)$; they are historical
outcomes of the actualization draw.
\item \textbf{Masses.} The Yukawa operator \emph{form} is derived; its \emph{spectrum} (the
hierarchy $1{:}207{:}3478$) is set by the contingent architecture shapes.
\item \textbf{Color count $n=3$.} The count is unfixed by any tested principle (provisional
gauge-count no-go T-09, 0.10); it is a \drawn\ value.
\end{itemize}

\subsection{The no-go theorems and the taxonomy at a glance}

The taxonomy classifies \emph{underivability}. The fourteen classified results, condensed,
are: $U(1)$ \derived\ (0.95); spatial 3 \derived\ (0.85); $N\ge3$ \derived\ (0.90);
the log-normal law \derived\ (theorem); $SU(2)$ \realunder\ (0.70); $SU(3)$ structure
\realunder\ (0.10); Koide $Q=2/3$ reality \realunder\ (0.90); Koide $45^\circ$ origin
\realunder\ (0.70); Yukawas/couplings/$\Omega_{DM}$ \drawn; $N\le3$ \drawn\ (0.70);
color count 3 \drawn; internal $N=3$ \derived$\cap$\drawn\ (0.70); and neutrino-Koide
\textbf{FALSIFIED} (0.90). Category conflicts: 0; composite objects: 2; overall
consistency: 0.95.

Five conditional no-go theorems delimit the boundary between the derivable and the
contingent: \textbf{T-08} (no symmetry/attractor/topology/information-geometry selects
$Q=2/3$; 0.70); \textbf{T-09} (no principle fixes the gauge count $n$; provisional, 0.10);
\textbf{T-10} (no stability/anomaly/representation bound gives $N\le3$; 0.70); \textbf{T-11}
(neutrino-Koide falsified; computation, 0.90); and \textbf{T-12} (one vs.\ three cascades
untestable from one universe; 0.55). The ``3'' that pervades physics collapses to one node
--- the \emph{internal-3} node --- with two formally unlinked faces ($N$ and $n$), and it is
the theory's one open door.

%=============================================================================
\section{The Continuum Program}
%=============================================================================

The continuum program establishes two controlled limits, and records that they are
disjoint.

\begin{itemize}
\item \textbf{$L_Q\to-\nabla^2$.} The chain Laplacian has the exact spectrum
$\lambda_k=(1/\Delta x^2)[2-2\cos(\pi k/(N{+}1))]$, converging to $(\pi k)^2$ with
second-order convergence (error $\sim4\times$/doubling). This is the elliptic, Riemannian
chain behind the Schr\"odinger program.
\item \textbf{BDG $\to\square$.} The causal-set Benincasa--Dowker--Glaser operator
converges to the d'Alembertian with leading error $O(h^2)$ (error $\sim4\times$/halving).
This is the Lorentzian, hyperbolic chain.
\item \textbf{The metric-origin chain.}
\begin{equation}
Q\text{-events}\to\text{causal order (native)}\to\text{conformal class (proven)}
\to\text{conformal factor (native)}\to g_{\mu\nu}.
\end{equation}
The conformal class is imported via Malament's theorem (a proven result, not a gap); the
conformal factor is native ($f=\rho^{2/d}$ from the counting measure). The metric
\emph{origin} is therefore closed, while the metric \emph{dynamics} remain imported.
\end{itemize}

The two chains are \emph{disjoint in signature} --- one elliptic and positive, the other
hyperbolic and indefinite --- so no native bridge between them has been closed. This is the
single most important structural fact of the continuum program.

%=============================================================================
\section{Verification}
%=============================================================================

All claims marked \textbf{verified} are backed by executable xUnit tests in
\texttt{TQM.Tests/ResearchXC} (or \texttt{ResearchQG}). The complete inventory is
\textbf{15 test files / 47 tests, all PASSING}:

\begin{center}
\begin{tabular}{lll}
\toprule
Chain link & Test file & Tests \\
\midrule
$L_Q\to-\nabla^2$ & GraphLaplacianContinuumTests & 1 \\
BDG$\to\square$ & BDGOperatorContinuumTests & 1 \\
$L_Q$ vs.\ BDG signature & QuantumGravityBridgeTests & 3 \\
no native $\Delta_g$ & CurvedSpaceBridgeTests & 3 \\
candidate constructions & MetricOperatorTests & 4 \\
weighted Laplacian & WeightedLaplacianTests & 4 \\
$L_W\to\Delta_g$ & LaplaceBeltramiTests & 3 \\
curved Schr\"odinger & CurvedSchrodingerTests & 3 \\
Einstein recovery & EinsteinRecoveryTests (PARTIAL) & 3 \\
Einstein chain & EinsteinTensorTests & 4 \\
Einstein integration & EinsteinTensorIntegrationTests & 4 \\
metric generation & MetricGenerationTests & 4 \\
metric emergence & MetricEmergenceTests & 4 \\
conformal structure & ConformalStructureTests & 3 \\
metric origin & MetricOriginTests & 3 \\
\bottomrule
\end{tabular}
\end{center}

\textbf{Decisive findings.} (1) The \emph{Schr\"odinger} side of the continuum limit is
controlled and tested: $L_Q\to-\nabla^2\to$ Schr\"odinger, plus the curved $L_W$, the
Laplace--Beltrami approximation, and unitarity. (2) The \emph{Einstein} side is partially
verified: the standard $g\to G_{\mu\nu}$ chain works, but TQM has no native metric (the
Malament import) and the BDG action is imported. (3) The conformal-class ``gap'' is an
imported proven theorem, not a theory gap.

\textbf{Method in brief.} Every substantive claim is made executable through an
\emph{audit}: a reproducible analyzer plus tests, with a disposition at the end
(\derived, \realunder, \drawn, no-go, blocked, contingent, selected). The no-go theorems
are produced by \emph{route exhaustion} under a no-new-primitives constraint, and are
therefore conditional, not absolute (except the computation T-11). Confidence numbers are
\emph{uncalibrated ordinal ranks}, not probabilities. The hostile-review history was almost
entirely a history of framing corrections --- ``closed'' $\to$ ``dispositioned'',
``no-go'' $\to$ ``conditional no-go'', ``derivation'' $\to$ ``ontological
reinterpretation'' --- rather than retracted physics; the Round-2 tally was 6 RESOLVED,
6 PARTIALLY RESOLVED, 0 OPEN.

%=============================================================================
\section{Remaining Open Issues}
%=============================================================================

Three scientific blockers remain, and they affect only a \emph{derivation-paper} claim:

\begin{enumerate}
\item \textbf{The G4 metric $\to$ operator coupling.} TQM imports $g_{\mu\nu}$ via
Malament rather than generating it from Q-events; the metric-dependent operator
$\Delta_g/\Box_g$ is absent. The metric \emph{origin} is closed, the metric
\emph{dynamics} are not.
\item \textbf{Native BDG derivation.} The Einstein--Hilbert side flows through the
imported causal-set BDG action; a native re-derivation from the Q-event primitives is
missing.
\item \textbf{A discriminating prediction.} No unique, sharp, currently-testable prediction
yet separates TQM from SM$+\Lambda$CDM. The RAR $g_\dagger=cH_0/2\pi$ matches $a_0$ but its
$2\pi$ factor is admitted accidental; $w(z)=-1+0.015(1+z)^{3/2}$ is a small, not-yet-detected
deviation.
\end{enumerate}

Additionally, several scope items are disclosed as admitted boundaries: the gauge-count
no-go T-09 is provisional (0.10); the CMB is an accepted partial computational layer; the
complexity argument is a window-intersection (T-02 $=$ 0.85); content is contingent by
design; the shared-cascade question is logically (not empirically) closed; and the unified
action is a roadmap.

%=============================================================================
\section{Publication Status}
%=============================================================================

\begin{center}
\fcolorbox{caveat}{white}{%
\begin{minipage}{0.94\textwidth}
\color{caveat}
\small
\centering
\textbf{READY AS A FOUNDATION MONOGRAPH}\\[1mm]
\textbf{READY AS A RESEARCH-PROGRAM PAPER}\\[1mm]
\textbf{NOT READY AS A DERIVATION PAPER}
\end{minipage}}
\end{center}

A foundation monograph may legitimately use imported proven theorems as inputs; the
remaining blockers affect only a derivation-paper claim. The Schr\"odinger dynamical core
is executable and tested; the Einstein recovery remains logical, not mathematical
(imported metric + imported BDG action, $G$ dimensional, no unique discriminating
prediction).

%=============================================================================
\section{Conclusion}
%=============================================================================

THE Q-MODEL is \emph{structurally complete} in the precise sense that every in-scope
question is \emph{dispositioned} --- derived, underivable, underdetermined, contingent, or
accepted as a computational layer. The single residual of genuine scientific tension is
the \emph{internal-3 node} (why the internal multiplicity/count saturates at 3), held open
by the provisional gauge-count no-go T-09 (0.10).

\begin{center}
\Large\bfseries Structure is derivable. Content is realized.
\end{center}

\end{document}
